% !TEX program = xelatex
% Generated by docs/latex/build.py from the sources pinned in sources.json.
\documentclass[11pt,a4paper]{article}
\usepackage[a4paper,margin=25mm]{geometry}
\usepackage{fontspec}
\setmainfont{lmroman10-regular.otf}[BoldFont=lmroman10-bold.otf,ItalicFont=lmroman10-italic.otf,BoldItalicFont=lmroman10-bolditalic.otf]
\setsansfont{lmsans10-regular.otf}[BoldFont=lmsans10-bold.otf,ItalicFont=lmsans10-oblique.otf]
\IfFontExistsTF{DejaVuSansMono.ttf}{\setmonofont{DejaVuSansMono.ttf}[Scale=0.83]}{\setmonofont{lmmono10-regular.otf}[Scale=0.83]}
\usepackage{amsmath,amssymb,mathtools}
\usepackage{graphicx,calc,longtable,booktabs,array}
\usepackage{fvextra}
\usepackage{xurl}
\usepackage[colorlinks=true,linkcolor=blue,urlcolor=blue,citecolor=blue]{hyperref}
\providecommand{\tightlist}{\setlength{\itemsep}{0pt}\setlength{\parskip}{0pt}}
\setlength{\parindent}{0pt}
\setlength{\parskip}{6pt}
\setlength{\emergencystretch}{3em}
\begin{document}
\section{P3 C3 --- Nontriangular Bulgarian solitaire: bound and all maximizers}\label{p3-c3-nontriangular-bulgarian-solitaire-bound-and-all-maximizers}

\href{https://hackathon.bainsa.ai/p/p3/c3}{Exact cell statement}. This is a written general proof. The accompanying Lean theorems certify finite ranks only; the finite replay below corroborates, but does not replace, the all-\(k\) argument.

A partition \(\lambda=(\lambda_1,\ldots,\lambda_s)\) represents \(n\) cards in positive piles. The Bulgarian move \(B\) subtracts one card from each pile, discards zero piles, adds a new pile of size \(s\), and sorts. A state is \emph{cyclic} if \(B^a(\lambda)=\lambda\) for some \(a>0\). Let \(d_B(\lambda)\) be the least \(t\ge0\) for which \(B^t(\lambda)\) is cyclic, and let \(D_B(n)=\max_{\lambda\vdash n}d_B(\lambda)\). Set \(T_j=j(j+1)/2\).

We use the established \href{https://github.com/mpelteshki/climbing-to-the-frontier/blob/212bba53b1818c08c40b8c4b48f636651eaf6410/cagent/p3/c1/written-proof.md}{C1 cyclic-boundary theorem}: at rank \(k\), where \(n=T_{k-1}+r\) and \(1\le r\le k\), the cyclic partitions are precisely the decreasing lists \(k-1+\varepsilon_0,k-2+\varepsilon_1,\ldots,\varepsilon_{k-1}\), omitting a final zero, with \(\varepsilon_j\in\{0,1\}\) and \(\sum\varepsilon_j=r\). Their boundary words rotate under \(B\), so their periods divide \(k\); each has pile sizes at most \(k\), and for \(1\le r<k\) the pile-count cycle includes a rise from \(k-1\) to \(k\). The cited C1 proof also establishes eventual cycling.

\subsection{Full C3 proof: nontriangular bound and extremal size}\label{full-c3-proof-nontriangular-bound-and-extremal-size}

Let \(T_j=j(j+1)/2\), and write every nontriangular card count as

\[
n=T_{k-1}+r,\qquad 1\le r\le k-1.
\]

For \(k\ge4\), the maximum first-cyclic depth satisfies

\[
D_B(n)\le k^2-2k-1.
\]

\begin{enumerate}
\def\labelenumi{(\Alph{enumi})}
\tightlist
\item
\end{enumerate}

At \(n=T_k-1\), equality holds. The section below also gives an exact, computable criterion for \textbf{every} maximizing starting partition at that size. This is a written proof; the general C3 statements have not been formalized in Lean. The existing Lean files verify the finite rank-4-through-rank-6 values. The published source is \href{https://sc.edu/study/colleges_schools/artsandsciences/mathematics/research/imi/research/documents/1998/1998_12.pdf}{Griggs and Ho, Theorem 4.4}; the argument here supplies the pile-lifetime details behind its upper bound and an explicit trajectory-count classification of equality.

\subsubsection{Sequence facts carried over from C2}\label{sequence-facts-carried-over-from-c2}

For a starting partition \(\lambda\), let \(c_i\) be the number of piles in \(B^{i-1}(\lambda)\). The new pile born on move \(i\) has size \(c_i\) and lives at times \(i+1,\ldots,i+c_i\). Denote its lifetime interval by \(J_i=[i+1,i+c_i]\). At time \(m\), exactly \(c_m\) intervals, including those of the original piles, are alive. Hence \(c_{i+1}\le c_i+1\); a rise by one means no pile dies on move \(i\).

The \href{https://github.com/mpelteshki/climbing-to-the-frontier/blob/212bba53b1818c08c40b8c4b48f636651eaf6410/cagent/p3/c2/upper-bound-proof.md}{C2 upper-bound proof} proves three sequence lemmas without using the C2 conclusion:

\begin{enumerate}
\def\labelenumi{\arabic{enumi}.}
\tightlist
\item
  If \(c_i\le x-1\) and \(c_j\ge x+1\) for \(i<j\), a \emph{sandwich pattern} \((c_p,\ldots,c_q)=(x-1,x,\ldots,x,x+1)\) occurs inside \([i,j]\).
\item
  Every such pattern of width \(L=q-p\ge2\) obeys \(p\le x(L-1)\). This follows from the explicit retreat to a shorter earlier pattern, stopping when \(p\le x\), as is forced at width two.
\item
  If \((c_p,\ldots,c_{p+k})=(k-2,k-1,\ldots,k-1,k)\), then \(p+k\le n+1\). The proof counts live intervals at time \(p+k\) and rules out an extra pile born after move 1.
\end{enumerate}

We also need two elementary links between a pile-count sequence and its state. First, its entire future determines its starting partition. At relative time \(m\), subtract from \(c_m\) the number of known new intervals \(J_i\) covering \(m\). The result is the number of original piles of size at least \(m\), for every \(m\), and these counts determine the partition. Thus, if \(c_{t+i}=c_{t+k+i}\) for all \(i\ge0\), the states at times \(t\) and \(t+k\) are equal; the state at time \(t\) is cyclic.

Second, let \(m\ge1\) satisfy

\[
c_m=k-1,\quad c_{m+1}=k,\quad
c_{m+a}\in\{k-1,k\}\ (2\le a<k),\quad c_{m+k}=k-1.
\]

\begin{enumerate}
\def\labelenumi{(\Alph{enumi})}
\setcounter{enumi}{1}
\tightlist
\item
\end{enumerate}

The state at time \(m\) is already cyclic. Indeed, for \(1\le a\le k\), all the \(a-1\) piles born on moves \(m,\ldots,m+a-2\) remain alive at time \(m+a-1\). If \(H_a\) counts the piles of the time-\(m\) state that have size at least \(a\), then

\[
H_a=c_{m+a-1}-(a-1)\in\{k-a,k-a+1\}.
\]

\begin{enumerate}
\def\labelenumi{(\Alph{enumi})}
\setcounter{enumi}{2}
\tightlist
\item
\end{enumerate}

At time \(m+k\), the \(k-1\) piles \(J_{m+1},\ldots,J_{m+k-1}\) already fill all \(c_{m+k}=k-1\) places, so no time-\(m\) pile has size above \(k\). Formula (C) says the conjugate diagram lies between the staircases of sizes \(k-1\) and \(k\); transposing gives the same containment for the original diagram. Its column heights therefore have the form \(k-1+\varepsilon_0,k-2+\varepsilon_1,\ldots,\varepsilon_{k-1}\), with bits \(\varepsilon_j\), so it is cyclic by the \href{https://github.com/mpelteshki/climbing-to-the-frontier/blob/212bba53b1818c08c40b8c4b48f636651eaf6410/cagent/p3/c1/written-proof.md}{C1 classification}.

\subsubsection{Finding a late sandwich}\label{finding-a-late-sandwich}

The C1 boundary classification also says every cyclic state at rank \(k\) has pile counts in \(\{k-1,k\}\) and period dividing \(k\). Because \(1\le r<k\), its boundary word contains both a zero and a one, so the cyclic count sequence has a rise \(k-1\) to \(k\). Choose the \textbf{smallest} positive index \(t\) for which

\[
c_t=k-1,\quad c_{t+1}=k,\quad
c_{t+i}=c_{t+k+i}\quad\text{for all }i\ge0.
\]

\begin{enumerate}
\def\labelenumi{(\Alph{enumi})}
\setcounter{enumi}{3}
\tightlist
\item
\end{enumerate}

This \(t\) exists. Let \(d_B(\lambda)\) be the first-cyclic depth. The future-count argument above gives \(d_B(\lambda)\le t-1\). If \(t\ge k+1\), minimality means the preceding block \((c_{t-k},\ldots,c_{t-1})\) differs from \((c_t,\ldots,c_{t+k-1})\); otherwise the periodicity and the rise in (D) would already begin at \(t-k\).

Assume \(t\ge k+1\). Since \(c_t=k-1\), one of the \(k\) recent intervals \(J_{t-k},\ldots,J_{t-1}\) is absent at time \(t\). Let \(u\) be the largest absent index. The rise \(c_t\to c_{t+1}\) entails no death, so the same lifetime calculation as in C2 gives

\[
c_u=t-u-1,\qquad c_{u+1}=t-u.
\]

\begin{enumerate}
\def\labelenumi{(\Alph{enumi})}
\setcounter{enumi}{4}
\tightlist
\item
\end{enumerate}

If \(u\ge t-k+1\), then \(c_u\le k-2\), and the sandwich rule between \(u\) and \(t+1\) produces a \textbf{type II} pattern

\[
(c_p,\ldots,c_q)=(k-2,k-1,\ldots,k-1,k),\quad
t-k+1\le p<q\le t+1,\quad q-p\le k.
\]

\begin{enumerate}
\def\labelenumi{(\Roman{enumi})}
\setcounter{enumi}{1}
\tightlist
\item
\end{enumerate}

Otherwise \(u=t-k\), so \(c_{t-k}=k-1\) and \(c_{t-k+1}=k\). If one of \(c_{t-k+2},\ldots,c_{t-1}\) is at most \(k-2\), the same rule gives (II). If one is at least \(k+1\), the rule between \(t-k\) and that index produces a \textbf{type I} pattern

\[
(c_p,\ldots,c_q)=(k-1,k,\ldots,k,k+1),\quad
t-k\le p<q\le t-1,\quad q-p\le k-1.
\]

\begin{enumerate}
\def\labelenumi{(\Roman{enumi})}
\tightlist
\item
\end{enumerate}

In the remaining case all intermediate counts lie in \(\{k-1,k\}\). Since (D) also gives \(c_t=k-1\), condition (B) holds at \(m=t-k\), so that earlier state is cyclic. Its future counts must then be \(k\)-periodic, contrary to the minimality of \(t\). Thus (I) or (II) always occurs.

\subsubsection{The upper bound}\label{the-upper-bound}

For \(k=4\), the exact \href{https://github.com/mpelteshki/climbing-to-the-frontier/blob/212bba53b1818c08c40b8c4b48f636651eaf6410/cagent/p3/c3/README.md}{Lean checks for \(n=7,8,9\)} give depths \(4,5,7\), all at most \(4^2-2\cdot4-1=7\). Let \(k\ge5\), and put \(B=k^2-2k-1\). If \(t\le k\), then \(d_B(\lambda)\le t-1\le k-1<B\). Otherwise take the pattern just found.

For type II of full width \(k\), its index range forces \((p,q)=(t-k+1,t+1)\). The special sequence lemma gives

\[
d_B(\lambda)\le t-1=p+k-2\le n-1<B.
\]

The last inequality follows from \(n\le T_k-1\) and \(T_k-2<B\) for \(k\ge5\).

Width \(k-1\) is impossible for either type. For type I, the index range forces \(q=t-1\), so \(c_{t-1}=k+1\). Move \(t-1\) then creates a pile of size \(k+1\) in the cyclic state at time \(t\), impossible because every rank-\(k\) cyclic pile has size at most \(k\). For type II, the index range and \(c_t=k-1\) force \((p,q)=(t-k+2,t+1)\): the only other possible placement would have \(q=t\) and \(c_t=k\). But \(c_p=k-2\) makes \(J_p\) end at \(p+k-2=t\), whereas \(c_t=k-1\to c_{t+1}=k\) permits no pile to end at \(t\).

Every remaining pattern has width \(L\le k-2\), level \(x\in\{k-1,k\}\), and \(p\ge t-k\). The C2 retreat bound yields

\[
d_B(\lambda)\le t-1\le p+k-1
\le x(L-1)+k-1
\le k(k-3)+k-1=B.
\]

This proves (A).

\subsubsection{\texorpdfstring{Equality at \(T_k-1\)}{Equality at T\_k-1}}\label{equality-at-t_k-1}

For every \(k\ge4\), begin with the staircase \((k,k-1,\ldots,1)\), delete its two cells on diagonal \(k\) in columns \(0,1\), and add one cell on diagonal \(k+1\) in column \(k\). The resulting partition of \(T_k-1\) is

\[
\lambda^*=(k-1,k-2,k-2,k-3,\ldots,1,1).
\]

\begin{enumerate}
\def\labelenumi{(\Alph{enumi})}
\setcounter{enumi}{5}
\tightlist
\item
\end{enumerate}

This is the witness of \href{https://sc.edu/study/colleges_schools/artsandsciences/mathematics/research/imi/research/documents/1998/1998_12.pdf}{Griggs and Ho, Theorem 4.4(2)}. Put \(B=k^2-2k-1\). For \(0\le t<B\), define a diagram \(D_t\) by deleting the diagonal-\(k\) cells in columns

\[
p_t=t\bmod k,\qquad p'_t=(p_t+1)\bmod k,
\]

and adding the diagonal-\((k+1)\) cell in column \(q_t=(k+t)\bmod(k+1)\). Here a cell on diagonal \(w\) in column \(j\) has row \(w-j\). We check that every \(D_t\) is a Ferrers partition. If \(p_t\le k-2\), the only invalid added-column positions are \(q_t\in\{p_t,p_t+1,p_t+2\}\): they create a vertical gap or reverse adjacent heights. If \(p_t=k-1\), the holes wrap to columns \(k-1,0\), and the invalid positions are \(q_t\in\{k-1,k,0,1\}\).

Write \(t=ak+b\), \(0\le b<k\). For \(t<B\), we have \(a\le k-3\), excluding \((a,b)=(k-3,k-1)\), and

\[
q_t=\begin{cases}
b-a-1,&b\ge a+1,\\
k+b-a,&b\le a.
\end{cases}
\]

\begin{enumerate}
\def\labelenumi{(\Alph{enumi})}
\setcounter{enumi}{6}
\tightlist
\item
\end{enumerate}

If \(b\le k-2\), the first branch gives \(q_t<b\), while the second gives \(q_t\ge b+3\), so all three forbidden positions are avoided. If \(b=k-1\), then \(a\le k-4\) and \(q_t=k-a-2\in\{2,\ldots,k-2\}\), avoiding all four wrapped forbidden positions. Thus every \(D_t\) before \(B\) is Ferrers. Before sorting, a Bulgarian move rotates each diagonal's column coordinates by one, so \(B(D_t)=D_{t+1}\) whenever \(t+1<B\). Since \(D_0=\lambda^*\), induction gives \(B^t(\lambda^*)=D_t\) for all \(t<B\).

For the Ferrers check above, when the added column differs from both deleted columns, its heights are explicitly \(h_j=k-j-\mathbf1_{j=p_t}-\mathbf1_{j=p'_t}+\mathbf1_{j=q_t}\) for \(0\le j\le k\). Adding at a deleted column leaves a vertical gap. In all other cases, inspecting adjacent differences \(h_j-h_{j+1}\) gives exactly the forbidden positions stated above; no other pair of heights reverses.

At \(t=B-1\), the holes are in columns \(k-2,k-1\) and the added cell is in column \(0\), giving

\[
B^{B-1}(\lambda^*)=(k+1,k-1,k-2,\ldots,3,1).
\]

Its next unsorted move has column heights \((k-1,k,k-2,\ldots,2)\); sorting the first two yields \((k,k-1,k-2,\ldots,2)\), a rank-\(k\) cyclic boundary partition of \(T_k-1\). Every earlier \(D_t\) has an occupied cell on diagonal \(k+1\), while a rank-\(k\) cyclic partition has none, by C1. Hence \(d_B(\lambda^*)=B\), and the upper bound proves \(D_B(T_k-1)=B\).

\subsubsection{Which starting partitions attain the maximum?}\label{which-starting-partitions-attain-the-maximum}

There is an exact pile-count signature for all maximizers. For \(k\ge5\), set \(n=T_k-1\), \(B=k^2-2k-1\), and \(P=k(k-3)\). A partition \(\lambda\vdash n\) has depth \(B\) \textbf{if and only if}

\[
\boxed{(c_P,c_{P+1},\ldots,c_{P+k-2})
       =(k-1,k,\ldots,k,k+1).}
\]

\begin{enumerate}
\def\labelenumi{(\Alph{enumi})}
\setcounter{enumi}{7}
\tightlist
\item
\end{enumerate}

The right side has \(k-1\) terms: one \(k-1\), then \(k-3\) copies of \(k\), then one \(k+1\). For necessity, equality in the upper-bound chain forces type I (type II has \(t\le p+k-1\)), level \(x=k\), width \(L=k-2\), and start \(p=k(k-3)\), giving (H). For sufficiency, its last term is \(c_{B-1}=k+1\). Move \(B-1\) creates a pile of size \(k+1\) in \(B^{B-1}(\lambda)\), so that state is not cyclic. Once a trajectory enters a cycle it never leaves; hence \(d_B(\lambda)\ge B\), and (A) gives equality. For \(k=4\), direct enumeration of the 30 partitions of \(9\) gives depth counts \(4,3,6,7,4,3,2,1\) at depths \(0,1,\ldots,7\), respectively. The unique depth-seven partition is \((3,2,2,1,1)\), which also satisfies (H) with \(P=4\). The exact Lean certificate verifies the maximum depth; the identifying enumeration can be replayed from recurrence (I) and the four cyclic boundary partitions.

The signature in fact fixes a particular late state. Let \(q=P+k-2=B-1\). At time \(q\), the piles born on moves \(P,\ldots,q-1\) are all alive. The pile \(J_P\), born with size \(k-1\), now has size \(2\); those born on moves \(P+1,\ldots,q-1\), each with initial size \(k\), have sizes \(4,5,\ldots,k\). These are \(k-2\) piles containing \(2+4+5+\cdots+k=T_k-4\) cards. Since \(c_q=k+1\) and \(n=T_k-1\), the other three piles have three cards in total, so each has size \(1\). Therefore every maximizer satisfies

\[
B^{B-2}(\lambda)=Q_k:=(k,k-1,\ldots,4,2,1,1,1).
\]

\begin{enumerate}
\def\labelenumi{(\Alph{enumi})}
\setcounter{enumi}{9}
\tightlist
\item
\end{enumerate}

Conversely, \(Q_k\) has depth exactly two. It has \(k+1\) piles, so its first move produces

\[
P_k=(k+1,k-1,k-2,\ldots,3,1),
\]

whose next move is the cyclic boundary \((k,k-1,\ldots,2)\). Both \(Q_k\) and \(P_k\) are noncyclic: the former has too many piles for a rank-\(k\) cyclic partition, and the latter has a pile larger than \(k\). Thus (J) is another necessary and sufficient condition for depth \(B\).

Equation (J) gives a finite \textbf{inverse construction of all starting maximizers}, using only a fixed seed and an explicit rule on partitions. For a partition \(\mu\) with \(m\) parts, choose any distinct part value \(s\) of \(\mu\) such that \(s\ge m-1\). Remove one occurrence of \(s\), add \(1\) to each remaining part, append \(s-m+1\) parts of size \(1\), and sort; call the result \(R_s(\mu)\). It has exactly \(s\) parts and \(B(R_s(\mu))=\mu\). Conversely every predecessor of \(\mu\) arises this way, because its new pile must have size \(s\), one of the parts of \(\mu\).

Start with \(\mathcal S_0=\{Q_k\}\) and repeat

\[
\mathcal S_{j+1}=\bigcup_{\mu\in\mathcal S_j}
  \{R_s(\mu):s\text{ is a distinct part of }\mu,\ s\ge\ell(\mu)-1\},
  \qquad 0\le j<B-2.
\]

\begin{enumerate}
\def\labelenumi{(\Alph{enumi})}
\setcounter{enumi}{10}
\tightlist
\item
\end{enumerate}

Then \(\mathcal S_{B-2}\) is \textbf{exactly} the set of all maximizing partitions of \(T_k-1\). Every state in that final set reaches \(Q_k\) after \(B-2\) moves, so has depth \(B\); every maximizer belongs to the set by (J) and the exhaustive predecessor rule. The construction never tests whether a state is cyclic and does not use a depth oracle. It can merge duplicate branches as sets. For \(k=4\), it yields the sole partition \((3,2,2,1,1)\); for larger \(k\), the witness (F) is one member among many.

Criterion (H) is a compact membership test for the same set. Its counts can be computed directly from the initial parts by the finite recurrence

\[
c_m=|\{j:\lambda_j\ge m\}|
  +|\{1\le i<m:i+c_i\ge m\}|,\qquad m\ge1.
\]

\begin{enumerate}
\def\labelenumi{(\Roman{enumi})}
\tightlist
\item
\end{enumerate}

The inverse construction (K) and count recurrence (I) are exact constructive answers to which partitions attain the maximum. They do not give a closed-form number of maximizers or a single static Ferrers-inequality description; neither is needed to generate or recognize every one.

\subsubsection{Exhaustive finite replay}\label{exhaustive-finite-replay}

Run \texttt{python3\ cagent/p3/c3/check.py} from the repository root. It independently computes forward depths and constructs the complete inverse sets for \(k=4,5,6,7\), then compares the sets elementwise. \href{https://github.com/mpelteshki/climbing-to-the-frontier/blob/212bba53b1818c08c40b8c4b48f636651eaf6410/cagent/p3/c3/finite-check.json}{Recorded results} have respectively 1, 6, 34, and 175 maximizing partitions. This also supplies the finite \(k=4\) classification used above.

\subsection{Appendix: pile-count lemmas used above}\label{appendix-pile-count-lemmas-used-above}

The next three sections reproduce the full general sequence arguments from the published \href{https://github.com/mpelteshki/climbing-to-the-frontier/blob/212bba53b1818c08c40b8c4b48f636651eaf6410/cagent/p3/c2/upper-bound-proof.md}{C2 upper-bound proof}. C3 uses these lemmas, not the triangular C2 conclusion.

\subsubsection{Pile lifetimes and a sequence bound}\label{pile-lifetimes-and-a-sequence-bound}

Write \(c_i\) for the number of piles in \(B^{i-1}(\lambda)\), with \(i\ge1\). An original pile of size \(a\) exists at times \(1,\ldots,a\). The pile created by move \(i\) has size \(c_i\) and exists at times

\[
J_i=[i+1,i+c_i].
\]

Its last time is \(e_i=i+c_i\). At any time \(m\), precisely \(c_m\) of these original-pile and created-pile intervals contain \(m\). One new interval starts between times \(i\) and \(i+1\); if \(d_i\) intervals end at \(i\), then

\[
c_{i+1}=c_i+1-d_i.
\]

\begin{enumerate}
\def\labelenumi{(\arabic{enumi})}
\tightlist
\item
\end{enumerate}

In particular \(c_{i+1}\le c_i+1\); a rise by one means \textbf{no} pile dies, and a constant pair means \textbf{exactly one} pile dies.

We use a \emph{sandwich pattern} of level \(x\) and width \(L=q-p\ge2\):

\[
(c_p,c_{p+1},\ldots,c_q)=(x-1,x,\ldots,x,x+1).
\]

\begin{enumerate}
\def\labelenumi{(\arabic{enumi})}
\setcounter{enumi}{1}
\tightlist
\item
\end{enumerate}

The two rises occur at its ends. The following elementary sandwich rule will locate one. If \(i<j\), \(c_i\le x-1\), and \(c_j\ge x+1\), choose \(p\) as the last index before \(j\) with \(c_p\le x-1\), and \(q\) as the first subsequent index with \(c_q\ge x+1\). Equation (1) forces \(c_p=x-1\), \(c_{p+1}=\cdots=c_{q-1}=x\), and \(c_q=x+1\). Thus (2) occurs within \([i,j]\).

\subsubsection{The retreat lemma}\label{the-retreat-lemma}

Suppose (2) occurs and \(p>x\). Among the \(x\) recent created piles \(J_{p-x},\ldots,J_{p-1}\), at least one is absent at time \(p\), since \(c_p=x-1\). The last, \(J_{p-1}\), is present. Let \(u\in[p-x,p-2]\) be the largest index whose interval is absent. All of \(J_{u+1},\ldots,J_{p-1}\) are present at \(p\). Since \(c_{p+1}=c_p+1\), none dies at \(p\), so they remain present at \(p+1\).

Put \(h=p-u-1\) and \(y=h+1=p-u\le x\). The absence of \(J_u\) at \(p\) gives \(c_u\le h\). The presence of \(J_{u+1}\) at \(p+1\) gives \(c_{u+1}\ge h+1\). Applying \(c_{u+1}\le c_u+1\) shows

\[
c_u=h=y-1,\qquad c_{u+1}=h+1=y.
\]

\begin{enumerate}
\def\labelenumi{(\arabic{enumi})}
\setcounter{enumi}{2}
\tightlist
\item
\end{enumerate}

Moreover \(J_{u+1}\) ends at time \(u+1+y=p+1\).

If \(L=2\), the pattern also has \(c_{p+2}=c_{p+1}+1\), so no interval ends at \(p+1\), a contradiction. Hence every width-two pattern satisfies

\[
p\le x.
\]

\begin{enumerate}
\def\labelenumi{(\arabic{enumi})}
\setcounter{enumi}{3}
\tightlist
\item
\end{enumerate}

Now let \(L\ge3\). As long as no later value \(c_{u+s}\) has risen from \(y\) to \(y+1\), the consecutive intervals

\[
J_{u+s},J_{u+s+1},\ldots,J_{p+s-1}
\]

are all present at time \(p+s\), for \(1\le s\le L-1\). This holds for \(s=1\) by the choice of \(u\) and the lack of a death at \(p\). If \(c_{u+s}=y\), then \(J_{u+s}\) ends at \((u+s)+y=p+s\). For \(s\le L-2\), the pair \(c_{p+s}=c_{p+s+1}=x\) permits exactly one death. It must be this interval; all the others survive, and \(J_{p+s}\) is born. This proves the induction. At each stage the active \(J_{u+s}\) implies \(c_{u+s}\ge y\), while (1) and the preceding value \(y\) imply \(c_{u+s}\le y+1\). Hence it either stays at \(y\) or supplies the desired rise.

At \(s=L-1\), the transition from \(c_{q-1}=x\) to \(c_q=x+1\) permits no death. Thus \(J_{u+L-1}\) cannot have size \(y\), which would make it end at \(q-1\). A first rise occurs at some \(v\) with

\[
u+2\le v\le u+L-1,
\qquad(c_u,\ldots,c_v)=(y-1,y,\ldots,y,y+1).
\]

\begin{enumerate}
\def\labelenumi{(\arabic{enumi})}
\setcounter{enumi}{4}
\tightlist
\item
\end{enumerate}

It occurs by \(p+1\): otherwise both \(c_p\) and \(c_{p+1}\) would equal \(y\), contradicting their values \(x-1,x\). Thus \(u\ge p-x\), \(v\le p+1\), \(y\le x\), and the new width \(v-u<L\). This is the retreat step.

Each retreat decreases the positive integer width. It therefore stops either at a pattern with start at most its level or at width two, where (4) gives the same bound. Working backward through \(p\le u+x\), while levels never increase, gives the explicit bound

\[
\boxed{p\le x(L-1)}
\]

\begin{enumerate}
\def\labelenumi{(\arabic{enumi})}
\setcounter{enumi}{5}
\tightlist
\item
\end{enumerate}

for every sandwich pattern (2). More formally, if \(p\le x\), then \(p\le x(L-1)\) since \(L\ge2\). Otherwise induction on \(L\) gives \(p\le u+x\le y(L'-1)+x\le x(L-2)+x=x(L-1)\) for the smaller width \(L'<L\). This abstract induction is also kernel-checked as \href{https://github.com/mpelteshki/climbing-to-the-frontier/blob/212bba53b1818c08c40b8c4b48f636651eaf6410/cagent/p3/lean/ProofPursuit/P3/C2SandwichBound.lean}{\texttt{sandwich\_start\_bound}}; the lifetime argument above proves its application hypotheses in writing.

\subsubsection{A special full-width pattern}\label{a-special-full-width-pattern}

We also need one bound that uses the number \(n\) of cards. Suppose

\[
(c_p,\ldots,c_{p+k})=(k-2,k-1,\ldots,k-1,k).
\]

\begin{enumerate}
\def\labelenumi{(\arabic{enumi})}
\setcounter{enumi}{6}
\tightlist
\item
\end{enumerate}

Here \(k\ge3\). The interval \(J_p\) ends at \(p+k-2\). Since \(c_{p+k-2}=c_{p+k-1}=k-1\), it is the \textbf{only} interval ending then. No interval ends at \(p+k-1\), since the next count rises from \(k-1\) to \(k\).

At time \(p+k\), the \(k-1\) intervals \(J_{p+1},\ldots,J_{p+k-1}\) are all present. There is exactly one other pile. If it is original, its original size is at least \(p+k\), and hence \(n\ge p+k\). If it is \(J_1\), then \(1+c_1\ge p+k\), and \(c_1\le n\) gives \(n+1\ge p+k\).

The only other possibility would be \(J_i\) for some \(2\le i\le p-1\); \(J_p\) has already died. Such a \(J_i\) is also the unique extra pile at time \(p+k-1\), alongside \(J_{p+1},\ldots,J_{p+k-2}\). Hence \(J_{i-1}\) is absent at that time, so \((i-1)+c_{i-1}\le p+k-2\). But \(J_i\) survives to \(p+k\), and (1) yields

\[
(i-1)+c_{i-1}\ge(i-1)+(c_i-1)\ge p+k-2.
\]

Equality follows, making \(J_{i-1}\) a second interval ending at \(p+k-2\), contrary to the uniqueness of \(J_p\). We have proved

\[
\boxed{p+k\le n+1}.
\]

\begin{enumerate}
\def\labelenumi{(\arabic{enumi})}
\setcounter{enumi}{7}
\tightlist
\item
\end{enumerate}
\end{document}
